\documentclass[a4paper, serif, 10pt]{moderncv}

%% moderncv
% color and style
\moderncvcolor{black}
\moderncvstyle{banking}

%% page layout/margins
\usepackage[top=2.5cm, bottom=2.5cm, left=1.5cm, right=1.5cm]{geometry}

\nopagenumbers{}

%% Babel config for Brazilian Portuguese language
% ref:
% - https://latex3.github.io/babel/guides/locale-portuguese.html
\usepackage[brazilian]{babel}

%% fontspec
\usepackage{fontspec}

\IfFontExistsTF{Linux Libertine}{
  \setmainfont[Ligatures={TeX, Common}, Numbers=OldStyle]{Linux Libertine}
}{}
\IfFontExistsTF{Linux Libertine O}{
  \setmainfont[Ligatures={TeX, Common}, Numbers=OldStyle]{Linux Libertine O}
}{}

%% CTeX
% CJK support, used to show CJK characters in the resume
%
% - heading=false: disable ctex heading
% - punct=kaiming: use kaiming punctuations styles for CJK
% - scheme=plain: use plain scheme, do not override \normalsize font size
% - space=auto: space settings for CJK characters
%
% CJK fonts are resolved per locale language, e.g. Noto Serif/Sans CJK JP
% for Japanese resumes, with fallback to the next region if not installed
%
% ref:
% - http://ctan.mirrorcatalogs.com/language/chinese/ctex/ctex.pdf
\usepackage[UTF8, heading=false, punct=kaiming, scheme=plain, space=auto]{ctex}

\IfFontExistsTF{Noto Serif CJK SC}{
  \setCJKmainfont{Noto Serif CJK SC}
}{}
\IfFontExistsTF{Noto Sans CJK SC}{
  \setCJKsansfont{Noto Sans CJK SC}
}{}

%% line spacing
\usepackage{setspace}
\setstretch{1.125}

%% URL styling - use normal text instead of monospace
\urlstyle{same}

%% Patch \httplink and \httpslink to handle full URLs
% The original moderncv macros always prepend http:// or https:// to URLs.
% This causes issues when the URL already has a protocol (e.g., https://example.com)
% resulting in malformed URLs like https://https://example.com.
% This patch checks if the URL already contains "://" and uses it directly if so.
% Note: We use \str_set:Nx (x-expansion) to expand macros like \@homepage first.
\ExplSyntaxOn
\str_new:N \g_yamlresume_url_str

\RenewDocumentCommand{\httpslink}{O{}m}{
  \str_set:Nx \g_yamlresume_url_str {#2}
  \str_if_in:NnTF \g_yamlresume_url_str {://}
    { \url{#2} }
    { \url{https://#2} }
}

\RenewDocumentCommand{\httplink}{O{}m}{
  \str_set:Nx \g_yamlresume_url_str {#2}
  \str_if_in:NnTF \g_yamlresume_url_str {://}
    { \url{#2} }
    { \url{http://#2} }
}
\ExplSyntaxOff

\name{Mariana Costa Almeida}{}
\title{Engenheira de QA | Automação de Testes | Qualidade de Software}
\phone[mobile]{+55 11 98765-4321}
\email{mariana.almeida@example.com}
\homepage{https://marianaalmeida.dev}

\address{Rua das Palmeiras, 123, São Paulo, SP, Brasil, 01234-567}{}{}

\extrainfo{{\small \faLinkedin}\ \href{https://www.linkedin.com/in/marianacalmeida}{@marianacalmeida} {} {} {} • {} {} {} 
{\small \faGithub}\ \href{https://github.com/marianaalmeida}{@marianaalmeida}}

\begin{document}

\maketitle

\section{Informações Básicas}

\cvline{}{Engenheira de QA com mais de 8 anos de experiência em garantia de qualidade de software, especializada em automação de testes e melhoria contínua de processos. Atuação em ambientes ágeis, colaborando com equipes multidisciplinares para entregar produtos com alta qualidade e confiabilidade.
%
\begin{itemize}
\item Experiência sólida em testes manuais e automatizados (web, API e mobile)
\item Domínio de ferramentas como Selenium, Cypress, Playwright, Appium e REST Assured
\item Conhecimento em CI/CD, Docker e boas práticas de DevOps
\item Forte habilidade em análise de requisitos, criação de casos de teste e documentação
\item Mentalidade analítica, proativa e orientada a resultados
\end{itemize}}
\section{Formação}

\cventry{fev. de 2015–dez. de 2019}
        {Bacharelado, Engenharia de Computação, Nota: 8.7}
        {Universidade de São Paulo}
        {\url{https://www.usp.br}}
        {}
        {\begin{itemize}
\item Formação sólida em engenharia de software, com ênfase em qualidade e testes
\item Participação em projetos de extensão voltados à inclusão digital
\end{itemize}
\textbf{Cursos}: Engenharia de Software, Banco de Dados, Redes de Computadores, Sistemas Operacionais, Algoritmos e Estruturas de Dados, Qualidade de Software, Testes de Software}
\section{Experiência Profissional}

\cventry{mar. de 2022–Atual}
        {Engenheira de QA Sênior}
        {InovaTech Solutions}
        {\url{https://www.inovatechsolutions.com.br}}
        {}
        {\begin{itemize}
\item Liderança técnica da equipe de QA, com 4 pessoas sob minha orientação
\item Definição e implementação da estratégia de testes automatizados para aplicações web e APIs
\item Criação de frameworks de automação com Cypress e Playwright, reduzindo o tempo de regressão em 60\%
\item Integração dos testes automatizados ao pipeline de CI/CD (GitLab CI), garantindo feedback rápido
\item Colaboração com squads de desenvolvimento e produto para garantir critérios de aceite claros e testáveis
\item Condução de análises de causa raiz e proposição de melhorias contínuas no processo de qualidade
\end{itemize}
\textbf{Palavras-chave}: Cypress, Playwright, CI/CD, Liderança, Qualidade}

\cventry{jan. de 2020–fev. de 2022}
        {Engenheira de QA Pleno}
        {DataTrust Systems}
        {\url{https://www.datatrustsystems.com}}
        {}
        {\begin{itemize}
\item Responsável pela garantia de qualidade de uma plataforma SaaS de gestão de dados
\item Desenvolvimento de testes automatizados de API utilizando REST Assured e Postman
\item Execução de testes manuais exploratórios e funcionais em ambientes web e mobile
\item Criação e manutenção de planos de teste, casos de teste e relatórios de bugs no Jira
\item Participação ativa em cerimônias ágeis (Scrum) e revisões de sprint
\item Apoio na implementação de testes de performance com JMeter
\end{itemize}
\textbf{Palavras-chave}: REST Assured, Postman, Jira, Scrum, JMeter}

\cventry{mar. de 2018–dez. de 2019}
        {Analista de QA Júnior}
        {QualiSoft Consultoria}
        {\url{https://www.qualisoft.com.br}}
        {}
        {\begin{itemize}
\item Execução de testes manuais em sistemas web e mobile para clientes de diversos setores
\item Elaboração de casos de teste e documentação de evidências
\item Reporte e acompanhamento de defeitos até a resolução
\item Suporte à automação de testes com Selenium WebDriver
\item Participação em treinamentos internos sobre boas práticas de qualidade
\end{itemize}
\textbf{Palavras-chave}: Selenium, Testes manuais, Documentação, Bug tracking}
\section{Idiomas}

\cvline{Português}{Nativo ou Bilíngue}
\cvline{Inglês}{Proficiência Profissional Fluente \hfill \textbf{Palavras-chave}: TOEFL 105, Leitura e escrita técnica}
\cvline{Espanhol}{Proficiência Profissional Limitada}
\section{Competências}

\cvline{Automação de Testes}{Especialista \hfill \textbf{Palavras-chave}: Selenium, Cypress, Playwright, Appium, REST Assured}
\cvline{Testes Manuais}{Avançado \hfill \textbf{Palavras-chave}: Testes exploratórios, Casos de teste, Testes de regressão, Testes de aceitação}
\cvline{Ferramentas de Gestão}{Avançado \hfill \textbf{Palavras-chave}: Jira, TestRail, Azure DevOps, Confluence}
\cvline{CI/CD e DevOps}{Intermediário \hfill \textbf{Palavras-chave}: GitLab CI, Jenkins, GitHub Actions, Docker}
\cvline{Programação}{Avançado \hfill \textbf{Palavras-chave}: Python, JavaScript, Java, SQL}
\cvline{Metodologias Ágeis}{Avançado \hfill \textbf{Palavras-chave}: Scrum, Kanban, BDD, TDD}
\section{Prêmios}

\cventry{dez. de 2023}
        {InovaTech Solutions}
        {Prêmio de Excelência em Qualidade}
        {}
        {}
        {Reconhecimento por liderar a iniciativa de automação de testes que reduziu em 60\% o tempo de regressão e aumentou a cobertura de testes para 85\%.}
\section{Certificados}

\cventry{jun. de 2020}
        {ISTQB}
        {ISTQB Certified Tester Foundation Level}
        {\url{https://www.istqb.org}}
        {}
        {}

\cventry{nov. de 2021}
        {ISTQB}
        {Certified Agile Tester}
        {\url{https://www.istqb.org}}
        {}
        {}
\section{Publicações}

\cventry{ago. de 2023}
        {Automação de Testes com Cypress: Boas Práticas}
        {Comunidade QA Brasil}
        {\url{https://medium.com/@marianaalmeida/automacao-de-testes-com-cypress}}
        {}
        {Artigo técnico abordando estratégias para criação de testes automatizados robustos e maintíveis com Cypress, incluindo organização de pastas, uso de fixtures e integração com CI/CD.}
\section{Projetos}

\cventry{jan. de 2023–Atual}
        {Framework modular para automação de testes web e API com Cypress e REST Assured}
        {Framework de Automação Open Source}
        {\url{https://github.com/marianaalmeida/qa-automation-framework}}
        {}
        {\begin{itemize}
\item Projeto open source para compartilhar boas práticas de automação de testes
\item Suporte a testes data-driven e geração de relatórios
\item Documentação completa e exemplos de uso
\end{itemize}
\textbf{Palavras-chave}: Cypress, REST Assured, Open Source, Automação}
\section{Interesses}

\cvline{Música}{Violão, Canto}
\cvline{Esportes}{Corrida, Yoga, Trilhas}
\section{Voluntariado}

\cventry{jan. de 2022–Atual}
        {Mentora}
        {Comunidade QA Brasil}
        {\url{https://www.comunidadeqabrasil.com.br}}
        {}
        {\begin{itemize}
\item Mentoria para profissionais iniciantes em QA, com foco em automação de testes e boas práticas
\item Participação em eventos e workshops da comunidade
\item Apoio na organização de meetups sobre qualidade de software
\end{itemize}}
    
\end{document}